% Generated draft: requires mathematical and source review.
\documentclass[11pt]{article}
\usepackage[margin=25mm]{geometry}
\usepackage{fontspec}
\setmainfont{Latin Modern Roman}
\usepackage{amsmath,amssymb,mathtools,unicode-math}
\setmathfont{Latin Modern Math}
\usepackage{enumitem,needspace,xurl,hyperref}
\setlength{\parindent}{0pt}
\setlength{\parskip}{5pt}
\setlength{\emergencystretch}{3em}
\newcommand{\N}{\mathbb{N}}
\newcommand{\Z}{\mathbb{Z}}
\newcommand{\Q}{\mathbb{Q}}
\newcommand{\R}{\mathbb{R}}
\newcommand{\C}{\mathbb{C}}
\newcommand{\F}{\mathbb{F}}
\newcommand{\A}{\mathbb{A}}
\newcommand{\PP}{\mathbb{P}}
\newcommand{\E}{\mathbb{E}}
\newcommand{\eps}{\varepsilon}
\newcommand{\dd}{\,\mathrm{d}}
\begin{document}
\section*{UnsolvedMath ID: 2309013}
\textbf{Source problem code:} catalogue entry
\subsection{Definitions and mathematical statement}
Let $K\subset\R^n$ be compact with $n\geq3$, and let $f:\R^n\to\R$ be continuous. Use $\mathcal D=C_c^\infty(\R^n)$ for test functions. For a nonzero test function $\phi$, write
\[
 d(\phi)=\operatorname{diam}(\operatorname{supp}\phi),\qquad
 \|\phi\|_*=\|\phi\|_\infty+d(\phi)\|\nabla\phi\|_\infty.
\]
The source calls $D(\phi)$ a least-diameter disc containing the support; in $\R^n$ the usual intended object is an enclosing Euclidean ball, but that terminology is not defined. It denotes by $C$ the harmonic capacity associated with the kernel $|x-y|^{2-n}$, without specifying normalization or its convention for noncompact difference sets.

\textbf{Approximation condition.} There exist functions $h_j$, each harmonic on an open neighborhood $U_j\supset K$, such that
\[
 \sup_{x\in K}|h_j(x)-f(x)|\longrightarrow0.
\]
The neighborhoods may depend on $j$; no fixed common neighborhood is required by the statement.

\textbf{Recorded capacity condition.} The source displays
\[
 \left|\int_{\R^n}f\,\delta\phi\,dx\right|
 \leq\eta(d(\phi))\|\phi\|_*\,C(D(\phi)-X), \tag{1}
\]
and asks whether the approximation condition is equivalent to this condition for all continuous $f$. However, $\delta\phi$ is not defined, $X$ is not defined, and the modulus function is called $n$ in the prose but $\eta$ in the formula.

If $\delta\phi$ means $\Delta\phi$, $X$ means $K$, the difference means set subtraction, and the modulus is consistently $\eta$, the natural version is: there exists $\eta:(0,\infty)\to[0,\infty)$ with $\eta(s)\to0$ as $s\downarrow0$ (nondecreasing in $s$ under the source's monotonic wording), such that for every test function
\[
 |\langle\Delta f,\phi\rangle|
 =\left|\int_{\R^n}f\Delta\phi\,dx\right|
 \leq\eta(d(\phi))\|\phi\|_*\,C(D(\phi)\setminus K). \tag{2}
\]
This is a conditional interpretation of (1), not an authenticated repair. The zero test function contributes only a trivial inequality once empty-support conventions are fixed.

The source reports that approximation implies the capacity condition, and that the converse holds when $f\in C^2$; it calls the condition formally analogous to one in rational approximation theory. These claims refer to the intended undamaged formulation. The supplied summary says the continuous-function question was open in the 2018 edition and no definitive resolution was located.
\subsection{Short English statement}
Is uniform approximation on a compact subset of $\R^n$, $n\geq3$, by locally harmonic functions characterized by the stated capacity bound on the distributional Laplacian? The source claims one direction and the smooth converse, but its printed inequality contains an undefined operator, set, and mismatched modulus symbol.
\subsection{Sources}
\begin{itemize}
\item Local UnsolvedMath registry, permanent ID 2309013, source code catalogue entry, ``Research Problems in Function Theory --- Problem 9.13''; attribution: A. G. O'Farrell. Supplied status open, classification OPEN-TRIAGE, updated 2026-08-06.
\item Hayman, \emph{Research Problems in Function Theory} (2018), Problem 9.13, \url{https://arxiv.org/abs/1809.07200}. The record reports source-TeX extraction and access on 2026-07-29; the 2018 update was retained for status triage. Current status and rights require review; the record describes publicly available arXiv source TeX.
\item The implication, $C^2$ converse, and rational-approximation analogy are stated without separate references; no missing citations are invented.
\end{itemize}
\end{document}
